% =============================================================================
%  V. METHODS: PRIMARY PREDICTION MODEL
% =============================================================================
\begin{frame}{V. Methods: Primary Prediction Model}
The main model is a \textbf{dynamic discrete-time competing-risks hazard
model}.

\vspace{6pt}
\begin{itemize}
  \item It predicts the \textbf{risk of sepsis every hour} during a patient's
        ICU stay.
  \item It is built using \textbf{multinomial pooled logistic regression}.
  \item D'Agostino et al.~\cite{dagostino1990pooled} showed that, under short
        time intervals, it gives results similar to a \textbf{time-dependent
        Cox regression}.
  \item Unlike the standard Cox model, it \textbf{does not require the
        proportional hazards assumption}.
  \item It also accounts for patients who \textbf{leave the ICU or die before
        developing sepsis} (competing risks).
\end{itemize}

\vspace{8pt}
\begin{exampleblock}{Models used for comparison}
Gradient-Boosted Trees (GBT), Static Cox model, NEWS2, qSOFA, and SIRS.
\end{exampleblock}
\end{frame}


% =============================================================================
%  V. METHODS: VALIDATION STRATEGY
% =============================================================================
\begin{frame}{V. Methods: Validation Strategy}
\renewcommand{\arraystretch}{1.15}\footnotesize
\begin{center}
\begin{tabular}{p{2cm}p{6.4cm}p{4.8cm}}
\toprule
\textbf{Validation} & \textbf{What was done} & \textbf{Why?} \\
\midrule
\textbf{Internal}
  & Trained on patients from \textbf{2008--2016} and tested on
    \textbf{2017--2019} data~\cite{guo2022temporal}.
  & To evaluate the model on newer, unseen patients. \\[2pt]
\textbf{Random split}
  & Trained on 75\% of patients and tested on 25\%, chosen
    randomly~\cite{guo2022temporal}.
  & To quantify how much the data split inflates reported
    performance (H4). \\[2pt]
\textbf{Pre-treatment}
  & Used only data collected \textbf{before the first antibiotic or blood
    culture}~\cite{kamran2024evaluation}.
  & To check whether the model is predicting sepsis or simply detecting when
    doctors start treatment. \\[2pt]
\textbf{External}
  & Tested the same trained model on the
    \textbf{eICU-CRD}~\cite{pollard2018eicu} dataset without retraining.
  & To see whether the model works in other hospitals. \\
\bottomrule
\end{tabular}
\end{center}

\vspace{2pt}
{\small\textbf{Main evaluation metrics:} \textbf{Utility
Score}~\cite{reyna2020utility} and \textbf{Calibration}.\par}

\vspace{1pt}
{\small\textbf{AUROC was reported only for comparison because it may not
detect clinically important performance
failures~\cite{wang2025srpm}.}\par}
\end{frame}


% =============================================================================
%  V. METHODS: METHODOLOGICAL COMMITMENTS
% =============================================================================
\begin{frame}{V. Methods: Methodological Commitments}
\begin{itemize}\itemsep4pt
  \item \textbf{Pre-registered protocol:} The study plan was fixed before
        model development to reduce bias (following TRIPOD+AI
        guidelines~\cite{collins2024tripodai}).
  \item \textbf{Leakage prevention:} The model only uses information
        available at each prediction time; missing values are handled using
        past information only.
  \item \textbf{Patient-level separation:} The same patient never appears in
        both training and testing data.
  \item \textbf{Frozen external validation:} The model trained on MIMIC-IV
        is directly tested on eICU-CRD without retraining.
  \item \textbf{Reproducible pipeline:} The complete workflow uses organised
        scripts, shared configurations, and publicly available code for
        replication.
\end{itemize}
\end{frame}
